\documentclass{article}
\usepackage{sty/Layout}
\usepackage{sty/Commands}
\usepackage{siunitx}
\usepackage{multirow}
\usepackage{marginnote}
\usepackage{xcolor}
\usepackage{array}
\usepackage{colortbl}
\usepackage{needspace}

% Textbreite 15 cm (210 mm - 2 x 30 mm)
\geometry{left=30mm,right=30mm}

\sisetup{output-decimal-marker = {,}, separate-uncertainty = true}

% ===========================================================================
% Kopfzeile
% ===========================================================================
\fancyhead[L]{\textsf{\textbf{Messprotokoll}}}
\fancyhead[C]{}
\fancyhead[R]{\textsf{Versuch 42 -- W\"armekapazit\"at}}

% ===========================================================================
% Makros fuer den Vordruck
% ===========================================================================
% Spaltentypen: L = linksbuendig, C = zentriert, A = Auswertungsspalte (grau)
\newcolumntype{L}[1]{>{\raggedright\arraybackslash}p{#1}}
\newcolumntype{C}[1]{>{\centering\arraybackslash}p{#1}}
\newcolumntype{A}[1]{>{\columncolor{black!8}\centering\arraybackslash}p{#1}}
\setlength{\tabcolsep}{3pt}

% Zeilenhoehe: \zh = 8,5 mm (normale Messzeile), \zhs = 7 mm (kompakte Zeile)
\newcommand{\zh}{\rule[-2.5mm]{0pt}{8.5mm}}
\newcommand{\zhs}{\rule[-2mm]{0pt}{7mm}}

% Leerlinie zum Ausfuellen, optionale Breite
\newcommand{\bl}[1][3cm]{\underline{\hspace{#1}}}

% Ankreuzfeld
\newcommand{\ak}{\raisebox{-0.4ex}{\setlength{\fboxsep}{0pt}\framebox[3.2mm]{\rule{0pt}{3.2mm}}}\,}

% Kurzbeschreibung: was wird in diesem Abschnitt gemacht
\newcommand{\was}[1]{%
  \par\noindent\fbox{\parbox{\dimexpr\linewidth-2\fboxsep-2\fboxrule\relax}{%
    \small\textbf{Ziel:} #1}}\par}

% Zeile im Einstellungsblock: Bezeichnung links, Felder rechts (umbrechend)
\newcommand{\einst}[2]{%
  \par\noindent\parbox[t]{3.3cm}{\raggedright #1}%
  \parbox[t]{11.7cm}{\raggedright #2\rule[-1.5mm]{0pt}{6mm}}\par}

% n Schreiblinien ueber die volle Breite
\newcount\lncnt
\newcommand{\linien}[1]{%
  \par\begingroup\setlength{\parskip}{0pt}\lncnt=0
  \loop\ifnum\lncnt<#1\advance\lncnt by 1
    \noindent\rule{0pt}{8mm}\rule{\linewidth}{0.4pt}\par
  \repeat\endgroup\par}

% Beobachtungsfeld mit konkreter Frage
\newcommand{\beob}[2]{%
  \par\noindent\begin{minipage}{\linewidth}
    \textbf{Beobachtung} (#1):\linien{#2}
  \end{minipage}\par}

% Hinweis unter Tabellen mit Auswertungsspalten
\newcommand{\auswhinweis}{%
  \par\vspace{0.5mm}\noindent
  {\footnotesize Graue Spalten mit (A): erst bei der Auswertung ausf\"ullen.}\par}

% Tabellenblock, der nicht ueber einen Seitenumbruch laeuft
\newenvironment{tblock}{\par\noindent\begin{minipage}{\linewidth}}{\end{minipage}\par}

% Abschnitt nur beginnen, wenn noch Platz fuer den Anfang ist
\newcommand{\abschnitt}[1]{\Needspace{9cm}\section{#1}}


\begin{document}

\pagestyle{fancy}

% ===========================================================================
% Kopf
% ===========================================================================
\begin{center}
  {\LARGE\bfseries Versuch 42 -- W\"armekapazit\"at} \\[0.4cm]
  {\large Messprotokoll}
\end{center}

\vspace{0.3cm}

\begin{tabular}{ll}
  \textbf{Gruppe:}  & 8 \\[10pt]
  \textbf{Autoren:} & \bl[8cm] \\[10pt]
                    & \bl[8cm] \\[10pt]
  \textbf{Tutor:}   & \bl[8cm] \\[10pt]
  \textbf{Datum:}   & \bl[8cm] \\[10pt]
  \textbf{Uhrzeit:} & von \bl[2.5cm] \quad bis \bl[2.5cm] \\[10pt]
\end{tabular}

\was{Mit der Mischungsmethode im Wasserkalorimeter werden die spezifischen
  W\"armekapazit\"aten von Graphit, Aluminium und Blei bestimmt. Zuerst wird der Wasserwert
  $W$ des Kalorimeters aus einem Mischungsversuch mit warmem Wasser ermittelt. Danach werden
  die gro\ss{}en Probek\"orper in siedendem Wasser auf $T_1$ gebracht und im Kalorimeter die
  Mischungstemperatur gemessen (Bereich \SIrange{20}{100}{\celsius}). Zuletzt werden die
  kleinen Probek\"orper in fl\"ussigem Stickstoff abgek\"uhlt und aus der verdampften
  Stickstoffmenge die spezifische W\"arme im Bereich \SIrange{-195.8}{20}{\celsius} bestimmt.
  Aus dem Verh\"altnis beider Molw\"armen folgen die Debye-Temperaturen.}

% ===========================================================================
% 1 Verwendete Geraete
% ===========================================================================
\section{Verwendete Ger\"ate}

\begin{tblock}
  {\footnotesize
  \begin{tabular}{|L{3.7cm}|L{3.2cm}||L{3.7cm}|L{3.2cm}|}
    \hline
    \zhs\textbf{laut I Messaufbau} & \textbf{Typ / Inventar-Nr.}
      & \textbf{zus\"atzlich laut VI} & \textbf{Typ / Inventar-Nr.} \\
    \hline\hline
    \zhs Kalorimeter mit Deckel & & Wasserkocher & \\ \hline
    \zhs Magnetr\"uhrer & & VE-Wasser & \\ \hline
    \zhs R\"uhrfisch & & Gitter / Netz im Kalorimeter & \\ \hline
    \zhs Heizplatte mit Glasbecher & & Tuch zum Abtrocknen & \\ \hline
    \zhs Stativ mit Drahthaken & & Faden zum Eintauchen & \\ \hline
    \zhs 3 gro\ss{}e Versuchsk\"orper\newline (Graphit, Alu, Blei) & & Millimeterpapier / Tablet & \\ \hline
    \zhs 3 kleine Versuchsk\"orper & & Zange / Handschuhe & \\ \hline
    \zhs Dewargef\"a\ss{} (fl.\ N$_2$) & & gro\ss{}er Vorrats-Dewar & \\ \hline
    \zhs elektronisches Thermometer & & weiteres: & \\ \hline
    \zhs elektronische Waage & & weiteres: & \\ \hline
    \zhs Stoppuhr & & weiteres: & \\ \hline
    \zhs Schutzbrille, Schutzhandschuhe & & weiteres: & \\ \hline
  \end{tabular}}

  \vspace{2mm}
  Max-Funktion am Thermometer vorhanden \ak \quad
  Schutzbrille und Handschuhe vorhanden und benutzt \ak \quad
  Kalorimeter innen und au\ss{}en trocken \ak
\end{tblock}

\subsection*{Ablesegenauigkeiten}

\begin{tblock}
  {\small
  \begin{tabular}{|L{3.0cm}|C{1.3cm}|C{2.2cm}|C{2.3cm}|C{2.3cm}|C{2.2cm}|}
    \hline
    \zhs\textbf{Instrument} & \textbf{Gr\"o\ss{}e}
      & \textbf{Messbereich\newline eingestellt}
      & \textbf{Skalenteilung /\newline Aufl\"osung}
      & \textbf{Ablese-\newline unsicherheit}
      & \textbf{Hersteller-\newline angabe} \\
    \hline\hline
    \zh Digitalthermometer & $T$ & & & & \\ \hline
    \zh Waage (Kalorimeter) & $m$ & & & & \\ \hline
    \zh Waage (Dewar) & $m$ & & & & \\ \hline
    \zh Stoppuhr & $t$ & & & & \\ \hline
    \zh Heizplatte (Skala) & $T_\mathrm{Platte}$ & & & & \\ \hline
    \zh Raumtemp.:\ \bl[1.4cm] & $T_\mathrm{Raum}$ & & & & \\ \hline
  \end{tabular}}

  \vspace{2mm}
  Reaktionszeit des Thermometers (Anzeige stabil nach ca.): \bl[1.8cm] s \quad
  Waage tariert \ak \quad Waage nivelliert \ak
\end{tblock}

\beob{Ist der Thermometerfehler laut Hersteller statistisch oder systematisch (siehe
  V.3)? Wurde f\"ur alle W\"agungen dieselbe Waage benutzt?}{3}

% ===========================================================================
% 2 Skizze des Versuchsaufbaus
% ===========================================================================
\abschnitt{Skizze des Versuchsaufbaus}

\begin{tblock}
  \textbf{Aufbau A -- Mischungskalorimeter (Aufgaben 1, 3 und 4).}
  {\small Heizplatte bzw.\ Wasserkocher $\to$ Glasbecher mit siedendem Wasser, Probek\"orper
  h\"angt am Drahthaken des Stativs frei im Wasser (ber\"uhrt den Boden nicht) $\to$ Umsetzen
  auf das Gitter im Kalorimetergef\"a\ss{} $\to$ Deckel sofort schlie\ss{}en. Kalorimeter steht
  zentrisch auf dem Magnetr\"uhrer; Probek\"orper au\ss{}ermittig (z.B.\ links), Thermometer
  auf der gegen\"uberliegenden Seite, ohne den K\"orper zu ber\"uhren.}

  \vspace{1mm}
  \framebox[\linewidth]{\rule{0pt}{7.1cm}}
\end{tblock}

\begin{tblock}
  \textbf{Aufbau B -- Dewar mit fl\"ussigem Stickstoff auf der Waage (Aufgabe 5).}
  {\small Dewargef\"a\ss{} zu ca.\ $3/4$ mit fl.\ N$_2$ gef\"ullt $\to$ steht w\"ahrend der
  gesamten Messung auf der Waage ($m(t)$) $\to$ kleiner Probek\"orper h\"angt am Faden und
  wird langsam eingetaucht $\to$ Stoppuhr l\"auft vom Eintauchen bis zum Ende des Siedens.}

  \vspace{1mm}
  \framebox[\linewidth]{\rule{0pt}{7.1cm}}
\end{tblock}

\beob{Abweichungen vom Aufbau in Abb.~1, Lage von Thermometer, K\"orper und R\"uhrfisch im
  Kalorimeter, Art der Wasserheizung (Heizplatte oder Wasserkocher)}{3}

% ===========================================================================
% 3 Aufgabe 1
% ===========================================================================
\abschnitt{Aufgabe 1: Vorbereiten des Wasserkalorimeters}

\was{Der Magnetr\"uhrer wird vor Versuchsbeginn erprobt: Kalorimetergef\"a\ss{} zu drei
  Vierteln mit VE-Wasser f\"ullen, zentrisch auf den R\"uhrer stellen und die Position so
  optimieren, dass das charakteristische Ratschen des Magneten deutlich zu h\"oren ist. Es
  werden keine Zahlenwerte gemessen, aber die gew\"ahlte Position gilt f\"ur alle folgenden
  Messungen.}

\subsection*{Einstellungen}

\einst{F\"ullung:}{Gef\"a\ss{} zu drei Vierteln mit VE-Wasser gef\"ullt \ak \quad
  F\"ullh\"ohe ca.\ \bl[1.5cm] cm}
\einst{R\"uhrer:}{eingeschaltet \ak \quad Stufe \bl[1.5cm] \quad
  Ratschen deutlich h\"orbar \ak \quad auch mit geschlossenem Deckel \ak}
\einst{Position:}{zentrisch \ak \quad verschoben um ca.\ \bl[1.5cm] cm nach \bl[2.5cm]}
\einst{Thermometer:}{taucht ein, ber\"uhrt den Probek\"orperplatz nicht \ak \quad
  Eintauchtiefe \bl[1.6cm] cm}
\einst{Uhrzeit:}{Vorbereitung abgeschlossen um \bl[1.8cm] Uhr}

\beob{Bei welcher R\"uhrerstufe wird gut durchmischt, ohne dass der R\"uhrfisch springt?
  Wie klingt das Ratschen bei richtiger Position, und bleibt sie w\"ahrend des Versuchs
  erhalten?}{3}

% ===========================================================================
% 4 Aufgabe 2
% ===========================================================================
\abschnitt{Aufgabe 2: Massen aller Probek\"orper}

\was{Alle sechs Probek\"orper werden trocken gewogen. Die gro\ss{}en K\"orper geh\"oren zu
  Aufgabe~4 (siedendes Wasser), die kleinen zu Aufgabe~5 (fl\"ussiger Stickstoff). Die Massen
  $m_x$ gehen unmittelbar in $c_x$ ein.}

\subsection*{Einstellungen}

\einst{Waage:}{tariert \ak \quad Anzeige stabil abgewartet \ak \quad
  Unsicherheit $\Delta m =$ \bl[1.8cm] g (f\"ur alle W\"agungen)}
\einst{K\"orper:}{trocken \ak \quad Faden bzw.\ Haken \ak mitgewogen \ak nicht mitgewogen}

\begin{tblock}
  {\small
  \begin{tabular}{|L{2.6cm}|C{1.6cm}|C{2.2cm}|C{1.8cm}|L{3.2cm}|A{1.9cm}|}
    \hline
    \zhs\textbf{Probek\"orper} & \textbf{Gr\"o\ss{}e} & \textbf{$m_x$\newline [g]}
      & \textbf{$\Delta m_x$\newline [g]} & \textbf{Zustand / Kennzeichnung}
      & \textbf{$n = m/M$ (A)\newline [mol]} \\
    \hline\hline
    \zh Graphit    & gro\ss{} & & & & \\ \hline
    \zh Aluminium  & gro\ss{} & & & & \\ \hline
    \zh Blei       & gro\ss{} & & & & \\ \hline
    \zh Graphit    & klein & & & & \\ \hline
    \zh Aluminium  & klein & & & & \\ \hline
    \zh Blei       & klein & & & & \\ \hline
  \end{tabular}}
  \auswhinweis
\end{tblock}

\beob{Oberfl\"ache der K\"orper (rau, oxidiert, feucht)? H\"angt an einem K\"orper noch
  Wasser oder ein Draht? Wurden die K\"orper vor dem Wiegen abgetrocknet?}{2}

% ===========================================================================
% 5 Aufgabe 3
% ===========================================================================
\abschnitt{Aufgabe 3: Bestimmung des Wasserwertes $W$}

\was{Das leere, trockene Kalorimeter (mit Netz und R\"uhrfisch, ohne Deckel) wird gewogen und
  seine Anfangstemperatur $T_2$ als Zimmertemperatur gemessen. Wasser wird auf etwa
  \SI{50}{\celsius} erhitzt, seine stabile Temperatur als $T_1$ notiert. Beim Einf\"ullen
  (Zeitpunkt $t_0$) startet die Uhr; danach wird f\"unf Minuten lang alle \SI{30}{\second}
  die Temperatur gemessen, bis ein linearer Abfall vorliegt. Die Mischungstemperatur $T$ wird
  sp\"ater zeichnerisch auf $t_0$ extrapoliert.}

\subsection*{W\"agungen und Ausgangstemperaturen}

\einst{Kalorimeter leer:}{trocken \ak \ mit Netz \ak \ mit R\"uhrfisch \ak \ ohne Deckel \ak
  \quad $m_\mathrm{Kal} =$ \bl[2.2cm] $\pm$ \bl[1.4cm] g}
\einst{Zimmertemperatur:}{$T_2 =$ \bl[2cm] $\pm$ \bl[1.4cm] \si{\celsius} \quad
  um \bl[1.6cm] Uhr}
\einst{Warmes Wasser:}{Sollwert ca.\ \SI{50}{\celsius}; Heizplatte abgeschaltet um
  \bl[1.6cm] Uhr}
\einst{}{Anzeige unter R\"uhren konstant \ak \quad
  $T_1 =$ \bl[2cm] $\pm$ \bl[1.4cm] \si{\celsius}}
\einst{Einf\"ullen:}{Kalorimeter zur H\"alfte gef\"ullt \ak \quad Deckel sofort aufgelegt \ak
  \quad R\"uhrer ein \ak}
\einst{}{Zeitpunkt $t_0 =$ \bl[1.8cm] Uhr \quad (Uhr gestartet \ak)}
\einst{Kalorimeter voll:}{nach der Messung, ohne Deckel: $m_\mathrm{Kal+W} =$ \bl[2.2cm]
  $\pm$ \bl[1.4cm] g}

\subsection*{Tabelle 3.1: Temperaturverlauf, alle \SI{30}{\second} \"uber \SI{5}{\minute}}

\begin{tblock}
  Unsicherheit: $\Delta T =$ \bl[1.8cm] \si{\celsius} \quad
  (Zeitnullpunkt $t = 0$ ist der Einf\"ullzeitpunkt $t_0$)

  \vspace{1mm}
  \begin{minipage}[t]{0.49\linewidth}
    {\small
    \begin{tabular}{|C{0.7cm}|C{1.5cm}|C{1.7cm}|C{2.0cm}|}
      \hline
      \zhs\textbf{Nr.} & \textbf{$t$\newline [s]} & \textbf{Uhrzeit}
        & \textbf{$T$\newline [\si{\celsius}]} \\
      \hline\hline
      \zhs 1 & 0   & & \\ \hline
      \zhs 2 & 30  & & \\ \hline
      \zhs 3 & 60  & & \\ \hline
      \zhs 4 & 90  & & \\ \hline
      \zhs 5 & 120 & & \\ \hline
      \zhs 6 & 150 & & \\ \hline
      \zhs 7 & 180 & & \\ \hline
    \end{tabular}}
  \end{minipage}\hfill
  \begin{minipage}[t]{0.49\linewidth}
    {\small
    \begin{tabular}{|C{0.7cm}|C{1.5cm}|C{1.7cm}|C{2.0cm}|}
      \hline
      \zhs\textbf{Nr.} & \textbf{$t$\newline [s]} & \textbf{Uhrzeit}
        & \textbf{$T$\newline [\si{\celsius}]} \\
      \hline\hline
      \zhs 8  & 210 & & \\ \hline
      \zhs 9  & 240 & & \\ \hline
      \zhs 10 & 270 & & \\ \hline
      \zhs 11 & 300 & & \\ \hline
      \zhs 12 & \bl[0.9cm] & & \\ \hline
      \zhs 13 & \bl[0.9cm] & & \\ \hline
      \zhs 14 & \bl[0.9cm] & & \\ \hline
    \end{tabular}}
  \end{minipage}

  \vspace{1mm}
  {\footnotesize Zeilen 12--14: Reserve, falls der lineare Abfall nach \SI{5}{\minute} noch
  nicht klar erkennbar ist.}
\end{tblock}

\subsection*{Ergebnisfelder (erst bei der Auswertung)}

\begin{tblock}
  {\small
  \begin{tabular}{|L{4.4cm}|A{3.2cm}|A{3.2cm}|L{3.0cm}|}
    \hline
    \zhs\textbf{Gr\"o\ss{}e} & \textbf{Wert (A)} & \textbf{Unsicherheit (A)}
      & \textbf{Bemerkung} \\
    \hline\hline
    \zh $m_W = m_\mathrm{Kal+W} - m_\mathrm{Kal}$ & & & \\ \hline
    \zh $\bar{T}$ (zeichnerisch auf $t_0$) & & & Millimeterpapier \\ \hline
    \zh $W$ nach Gl.~(13) & & & \\ \hline
  \end{tabular}}
  \auswhinweis
\end{tblock}

\beob{Wie schnell f\"allt die Temperatur, ab welchem Zeitpunkt ist der Abfall linear? Wurde
  beim Einf\"ullen Wasser versch\"uttet oder ist Dampf entwichen? Lief der R\"uhrer w\"ahrend
  der gesamten Messung?}{3}

% ===========================================================================
% 6 Aufgabe 4
% ===========================================================================
\abschnitt{Aufgabe 4: Spezifische W\"armekapazit\"aten mit kochendem Wasser}

\was{Das Kalorimeter wird zu drei Vierteln mit frischem VE-Wasser gef\"ullt und gewogen. Der
  gro\ss{}e Probek\"orper h\"angt mindestens \SI{5}{\minute} im sprudelnd kochenden Wasser und
  nimmt dessen Temperatur $T_1$ an. Vor dem Einbringen wird die Kalorimetertemperatur einige
  Minuten lang jede Minute gemessen ($T_2$), nach dem Einbringen bis zum sicheren Erreichen
  des Temperaturmaximums $T$. Vor jeder Messung werden Wassermasse und $T_2$ neu bestimmt.}

\subsection*{Gemeinsame Einstellungen}

\einst{Heizplatte:}{beim Aufheizen h\"ochstens \SI{300}{\celsius} \ak \quad nach dem
  Aufkochen auf \SI{250}{\celsius} reduziert \ak \quad eingestellt \bl[1.6cm] \si{\celsius}}
\einst{Glasbecher:}{zur H\"alfte gef\"ullt \ak \quad ca.\ \bl[1.5cm] ml \quad
  Wasser sprudelnd \ak}
\einst{Probek\"orper:}{h\"angt am Stativ \ak \quad ber\"uhrt den Becherboden nicht \ak \quad
  vollst\"andig untergetaucht \ak}
\einst{Kalorimeter:}{zu drei Vierteln mit frischem VE-Wasser \ak \quad R\"uhrer ein \ak \quad
  Deckel sofort nach dem Einsetzen \ak}
\einst{Waage:}{$\Delta m =$ \bl[1.6cm] g \quad Thermometer: $\Delta T =$ \bl[1.6cm]
  \si{\celsius}}

% --- Makro fuer einen Probekoerper -----------------------------------------
\newcommand{\koerperblock}[2]{%
  \Needspace{12cm}
  \subsection*{#1}

  \begin{tblock}
    \einst{W\"agung:}{$m_\mathrm{Kal+W} =$ \bl[2.2cm] g \quad um \bl[1.5cm] Uhr \quad
      (leer: $m_\mathrm{Kal} =$ \bl[2.2cm] g)}
    \einst{Siedebad:}{$T_1 =$ \bl[2cm] $\pm$ \bl[1.3cm] \si{\celsius} \quad
      Anzeige stabil \ak \quad gemessen um \bl[1.5cm] Uhr}
    \einst{Verweildauer:}{von \bl[1.5cm] bis \bl[1.5cm] Uhr \quad
      (mindestens \SI{5}{\minute} \ak)}
    \einst{Einsetzen:}{Zeitpunkt $t_0 =$ \bl[1.6cm] Uhr \quad Umsetzdauer ca.\ \bl[1.2cm] s
      \quad Deckel sofort zu \ak}
  \end{tblock}

  \vspace{1mm}
  \begin{tblock}
    \begin{minipage}[t]{0.40\linewidth}
      \textbf{#2a Vorlauf (jede Minute)}\\[1mm]
      {\small
      \begin{tabular}{|C{0.7cm}|C{1.4cm}|C{1.9cm}|}
        \hline
        \zhs\textbf{Nr.} & \textbf{Uhrzeit} & \textbf{$T$\newline [\si{\celsius}]} \\
        \hline\hline
        \zhs 1 & & \\ \hline
        \zhs 2 & & \\ \hline
        \zhs 3 & & \\ \hline
        \zhs 4 & & \\ \hline
        \zhs 5 & & \\ \hline
      \end{tabular}}

      \vspace{1mm}
      {\footnotesize $T_2$ = letzter Wert vor dem Einsetzen.}

      \vspace{1mm}
      $T_2 =$ \bl[1.6cm] \si{\celsius}
    \end{minipage}\hfill
    \begin{minipage}[t]{0.58\linewidth}
      \textbf{#2b Verlauf nach dem Einsetzen (alle \SI{30}{\second})}\\[1mm]
      {\small
      \begin{tabular}{|C{0.7cm}|C{1.2cm}|C{1.4cm}|C{1.9cm}|C{1.5cm}|}
        \hline
        \zhs\textbf{Nr.} & \textbf{$t$\newline [s]} & \textbf{Uhrzeit}
          & \textbf{$T$\newline [\si{\celsius}]} & \textbf{steigt /\newline f\"allt} \\
        \hline\hline
        \zhs 1  & 0   & & & \\ \hline
        \zhs 2  & 30  & & & \\ \hline
        \zhs 3  & 60  & & & \\ \hline
        \zhs 4  & 90  & & & \\ \hline
        \zhs 5  & 120 & & & \\ \hline
        \zhs 6  & 150 & & & \\ \hline
        \zhs 7  & 180 & & & \\ \hline
        \zhs 8  & 210 & & & \\ \hline
        \zhs 9  & 240 & & & \\ \hline
        \zhs 10 & 270 & & & \\ \hline
        \zhs 11 & 300 & & & \\ \hline
        \zhs 12 & \bl[0.8cm] & & & \\ \hline
      \end{tabular}}
    \end{minipage}
  \end{tblock}

  \vspace{1mm}
  \begin{tblock}
    \einst{Maximum:}{$T =$ \bl[2cm] $\pm$ \bl[1.3cm] \si{\celsius} \quad
      um \bl[1.5cm] Uhr \quad \ak abgelesen \ \ak Max-Funktion}
    \einst{Kontrolle:}{Temperatur f\"allt danach wieder \ak \quad
      K\"orper ber\"uhrte das Thermometer nicht \ak}
  \end{tblock}
  \beob{Wie lange dauert es bis zum Maximum? Spritzt Wasser beim Einsetzen? H\"angt Wasser
    am K\"orper, oder ist er beim Umsetzen sichtbar abgek\"uhlt?}{2}}

\koerperblock{Tabelle 4.1: Graphit (gro\ss{})}{4.1}

\koerperblock{Tabelle 4.2: Aluminium (gro\ss{})}{4.2}

\koerperblock{Tabelle 4.3: Blei (gro\ss{})}{4.3}

\subsection*{Zusammenfassung Aufgabe 4}

\begin{tblock}
  {\small
  \begin{tabular}{|L{2.0cm}|C{1.5cm}|C{1.5cm}|C{1.5cm}|C{1.5cm}|A{1.6cm}|A{1.6cm}|A{1.7cm}|}
    \hline
    \zhs\textbf{Probe} & \textbf{$m_x$\newline [g]} & \textbf{$m_W$\newline [g]}
      & \textbf{$T_1$\newline [\si{\celsius}]} & \textbf{$T_2$\newline [\si{\celsius}]}
      & \textbf{$T$\newline [\si{\celsius}]}
      & \textbf{$c_x$ (A)\newline [J/(gK)]}
      & \textbf{$c_\mathrm{mol}$ (A)\newline [J/(mol\,K)]} \\
    \hline\hline
    \zh Graphit   & & & & & & & \\ \hline
    \zh Aluminium & & & & & & & \\ \hline
    \zh Blei      & & & & & & & \\ \hline
  \end{tabular}}

  \vspace{0.5mm}
  {\footnotesize $m_W$ und $T$ werden hier aus den Bl\"ocken 4.1--4.3 \"ubertragen;
  $T$ ist das gemessene Maximum, also ein Messwert.}
  \auswhinweis
\end{tblock}

% ===========================================================================
% 7 Aufgabe 5
% ===========================================================================
\abschnitt{Aufgabe 5: Spezifische W\"armekapazit\"aten mit fl\"ussigem Stickstoff}

\was{Der Dewar wird zu etwa $3/4$ mit fl\"ussigem Stickstoff gef\"ullt, auf die Waage
  gestellt und der Temperaturausgleich abgewartet. Aus zwei W\"agungen im Abstand von
  \SI{2}{\minute} folgt die Verdunstungsrate unter Umgebungsbedingungen. Danach wird
  unmittelbar vor dem Eintauchen erneut gewogen, der kleine Probek\"orper (Anfangstemperatur
  $T_1 = $ Zimmertemperatur) langsam am Faden eingetaucht und die Zeit bis zum Ende des
  Siedens gestoppt. Die Massendifferenz ergibt die verdampfte Stickstoffmenge $m_V$.}

\subsection*{Einstellungen}

\einst{Sicherheit:}{Schutzbrille \ak \quad Handschuhe \ak \quad Raum gel\"uftet \ak}
\einst{Dewar:}{zu ca.\ $3/4$ gef\"ullt \ak \quad steht frei auf der Waage \ak \quad
  Waage tariert \ak}
\einst{Temperaturausgleich:}{abgewartet von \bl[1.5cm] bis \bl[1.5cm] Uhr \quad
  Sieden abgeklungen \ak}
\einst{Zimmertemperatur:}{$T_1 =$ \bl[2cm] $\pm$ \bl[1.3cm] \si{\celsius} \quad
  um \bl[1.5cm] Uhr}
\einst{Waage:}{$\Delta m =$ \bl[1.6cm] g \quad Anzeige driftet: \ak nein \ \ak ja, ca.\
  \bl[1.4cm] g/min}
\einst{Eintauchen:}{langsam am Faden \ak \quad K\"orper vollst\"andig bedeckt \ak \quad
  Uhr sofort gestartet \ak}

\subsection*{Tabelle 5.1: Verdunstungsrate (W\"agung im Abstand von \SI{2}{\minute})}

\begin{tblock}
  {\small
  \begin{tabular}{|L{2.0cm}|C{1.9cm}|C{1.3cm}|C{1.9cm}|C{1.3cm}|C{1.5cm}|A{2.2cm}|}
    \hline
    \zhs\textbf{Durchgang} & \textbf{$m_\mathrm{A}$\newline [g]} & \textbf{Uhrzeit}
      & \textbf{$m_\mathrm{B}$\newline [g]} & \textbf{Uhrzeit}
      & \textbf{$\Delta t$\newline [s]}
      & \textbf{Rate (A)\newline [g/min]} \\
    \hline\hline
    \zh vor Graphit   & & & & & & \\ \hline
    \zh vor Aluminium & & & & & & \\ \hline
    \zh vor Blei      & & & & & & \\ \hline
  \end{tabular}}

  \vspace{0.5mm}
  {\footnotesize Die Rate wird vor jedem Probek\"orper neu bestimmt, weil sie vom F\"ullstand
  und von den Umgebungsbedingungen abh\"angt.}
  \auswhinweis
\end{tblock}

\subsection*{Tabelle 5.2: Abk\"uhlen der kleinen Probek\"orper}

\begin{tblock}
  {\small
  \begin{tabular}{|L{1.7cm}|C{1.8cm}|C{1.2cm}|C{1.2cm}|C{1.3cm}|C{1.8cm}|A{1.7cm}|A{1.9cm}|}
    \hline
    \zhs\textbf{Probe} & \textbf{$m_0$ vor dem\newline Eintauchen\newline [g]}
      & \textbf{Ein-\newline tauchen\newline [Uhr]}
      & \textbf{Siede-\newline ende\newline [Uhr]}
      & \textbf{$\Delta t$\newline [s]}
      & \textbf{$m_\mathrm{end}$\newline [g]}
      & \textbf{$m_V$ (A)\newline [g]}
      & \textbf{$c_x$ (A)\newline [J/(gK)]} \\
    \hline\hline
    \zh Graphit   & & & & & & & \\ \hline
    \zh Aluminium & & & & & & & \\ \hline
    \zh Blei      & & & & & & & \\ \hline
  \end{tabular}}

  \vspace{0.5mm}
  {\footnotesize $m_\mathrm{end}$ wird direkt nach dem Herausnehmen des K\"orpers abgelesen.
  $m_V$ ist die um die Verdunstungsrate aus Tabelle~5.1 korrigierte Massendifferenz.}
  \auswhinweis
\end{tblock}

\subsection*{Kontrollen}

\begin{tblock}
  \einst{Reif / Eis:}{am Faden \ak \quad am K\"orper beim Herausnehmen \ak \quad
    am Dewarrand \ak}
  \einst{Rest-Stickstoff:}{in den gro\ss{}en Dewar zur\"uckgesch\"uttet \ak \quad
    um \bl[1.6cm] Uhr}
  \einst{Ende der Messung:}{\bl[1.8cm] Uhr \quad Raumtemperatur am Ende
    \bl[1.8cm] \si{\celsius}}
\end{tblock}

\beob{Woran ist das Ende des Siedens zu erkennen, und wie eindeutig ist dieser Zeitpunkt?
  Wie stark schwankt die Waagenanzeige w\"ahrend des Siedens? Bildet sich Eis, das die
  W\"agung verf\"alscht?}{4}

% ===========================================================================
% 8 Nennwerte und Konstanten
% ===========================================================================
\abschnitt{Nennwerte und Konstanten laut Anleitung}

\begin{tblock}
  {\small
  \begin{tabular}{|L{6.0cm}|L{5.4cm}|L{2.6cm}|}
    \hline
    \zhs\textbf{Gr\"o\ss{}e} & \textbf{Wert} & \textbf{Quelle} \\
    \hline\hline
    \zhs Spez.\ W\"arme von Wasser (\SIrange{15}{65}{\celsius})
      & $c_W = (4{,}186 \pm 0{,}004)$ J/(g\,K) & V.2 \\ \hline
    \zhs Wasserwert des Kalorimeters (Herstellerangabe)
      & $W \approx \SI{70}{\joule\per\kelvin}$ & V.2 \\ \hline
    \zhs Siedetemperatur fl.\ Stickstoff & $T_S = \SI{-195.8}{\celsius}$ & Anhang VIII \\ \hline
    \zhs Verdampfungsw\"arme fl.\ Stickstoff & $Q_V = \SI{199}{\joule\per\gram}$ & Anhang VIII \\ \hline
    \zhs Rel.\ Atommasse Blei & $M_\mathrm{Pb} = 207{,}2$ g/mol & Anhang VIII \\ \hline
    \zhs Rel.\ Atommasse Aluminium & $M_\mathrm{Al} = 26{,}98$ g/mol & Anhang VIII \\ \hline
    \zhs Rel.\ Atommasse Graphit & $M_\mathrm{C} = 12{,}01$ g/mol & Anhang VIII \\ \hline
    \zhs Literaturwert $c$ Blei (\SI{300}{\kelvin}) & $0{,}129$ J/(g\,K) & Anhang VIII \\ \hline
    \zhs Literaturwert $c$ Aluminium (\SI{300}{\kelvin}) & $0{,}90$ J/(g\,K) & Anhang VIII \\ \hline
    \zhs Literaturwert $c$ Graphit (\SI{300}{\kelvin}) & $0{,}709$ J/(g\,K) & Anhang VIII \\ \hline
    \zhs Debye-Temperatur Blei & $\Theta_{D,\mathrm{Pb}} = \SI{95}{\kelvin}$ & Abb.~2 \\ \hline
    \zhs Debye-Temperatur Aluminium & $\Theta_{D,\mathrm{Al}} = \SI{430}{\kelvin}$ & Abb.~2 \\ \hline
    \zhs Debye-Temperatur Diamant & $\Theta_{D,\mathrm{C}} = \SI{2230}{\kelvin}$ & Abb.~2 \\ \hline
    \zhs Ideale Gaskonstante & $R = 8{,}314$ J/(mol\,K) & Gl.~(4) \\ \hline
  \end{tabular}}

  \vspace{1mm}
  {\footnotesize F\"ur Graphit gibt es keinen einheitlichen Debye-Temperaturwert; der
  Diamantwert darf nur f\"ur einen qualitativen Vergleich benutzt und muss diskutiert werden.}
\end{tblock}

\vspace{2mm}
\noindent\textbf{Zusammenh\"ange:}
\begin{align*}
  &C = \frac{Q}{\Delta T}, \qquad c = \frac{Q}{m\,\Delta T}, \qquad
    c_\mathrm{mol} = \frac{M}{m}\,\frac{Q}{\Delta T}, \qquad
    W = m_W c_W\,\frac{T_1 - \bar{T}}{\bar{T} - T_2}, \\
  &c_x = \frac{(m_W c_W + W)\,(T - T_2)}{m_x\,(T_1 - T)}, \qquad
    Q = Q_V\,m_V = m_x c_x\,(T_1 - T_2), \qquad
    c_\mathrm{mol,DP} = 3R.
\end{align*}

% ===========================================================================
% 9 Bemerkungen
% ===========================================================================
\abschnitt{Bemerkungen und Besonderheiten w\"ahrend der Messung}

\beob{St\"orungen, Umbauten, Abweichungen von der Anleitung, Ger\"atewechsel, versch\"uttetes
  Wasser, wer hat was abgelesen?}{8}

\vspace{1.5cm}

\noindent
\begin{tabular}{@{}p{6.8cm}@{\hspace{1.2cm}}p{6.8cm}@{}}
  \rule{6.8cm}{0.4pt} & \rule{6.8cm}{0.4pt} \\
  Datum, Unterschrift Studierende & Datum, Unterschrift Tutor:in \\
\end{tabular}

% ===========================================================================
% Zusatzseite
% ===========================================================================
\newpage
\section*{Zus\"atzlicher Platz f\"ur Messwerte und Notizen}

\begin{tblock}
  {\small
  \begin{tabular}{|C{0.6cm}|C{1.5cm}|C{2.4cm}|C{2.4cm}|C{2.4cm}|C{2.4cm}|C{1.8cm}|}
    \hline
    \zhs\textbf{Nr.} & \textbf{Uhrzeit} & & & & & \\
    \hline\hline
    \zh & & & & & & \\ \hline
    \zh & & & & & & \\ \hline
    \zh & & & & & & \\ \hline
    \zh & & & & & & \\ \hline
    \zh & & & & & & \\ \hline
    \zh & & & & & & \\ \hline
    \zh & & & & & & \\ \hline
    \zh & & & & & & \\ \hline
    \zh & & & & & & \\ \hline
    \zh & & & & & & \\ \hline
  \end{tabular}}

  \vspace{1mm}
  {\footnotesize Leere Kopfzeile: Gr\"o\ss{}e und Einheit selbst eintragen, Unsicherheit als
  eigene Spalte.}
\end{tblock}

\linien{12}

\end{document}
